%!TEX root = main.tex
\section{Analysis}
\label{sec:analysis}

\subsection{Provenance of the 38 Changes (ViHSD)}
\label{sec:provenance}

The aggregate result is neutral; the question is what the 38 changes consist of.
Of the 23 corrections, 16 are \hate{} recoveries (7 from \off{}, 9 from \clean{})
and 7 are \clean{} de-flagging cases (6 from \hate{}, 1 from \off{}).  All 16
\hate{} recoveries have a rule firing with a \hate{} score above the \off{}
score, i.e.\ the rule points in exactly the corrected direction.  The rule
\hate{} score has median 0.7652 (mean 0.7507, range 0.2351--0.9939).

All 7 \clean{} de-flagging cases have a silent rule ($r_H=r_O=0$, so
$\tilde{\mathbf r}=(0,0,1)$).  With $\alpha_H=0.2$ the rule term disappears from
the \hate{} channel and $f_H=0.2\,p_H$, so a comment needs $p_H\ge0.70$ merely
to reach $\tau_H=0.14$; one flagged \hate{} at $p_H\approx0.5$ ends at
$f_H\approx0.10$ and the decision falls back to \clean{}.

The 15 regressions split into 10 rule-driven and 5 threshold-driven cases, and
the split aligns perfectly with direction: all 10 rule-driven cases
(\off{}$\to$\hate{}, \clean{}$\to$\hate{}) have the rule firing in the \hate{}
direction, so fusion trusts lexical \hate{} evidence over the neural label and
the lexicon is wrong about intensity or target.  The \off{} cases are the
expected failure of a phrase lexicon: profanity used as an intensifier rather
than an attack.  The 5 threshold-driven regressions are the mirror image: the
rule is silent, the fused \hate{} score is $0.2\,p_H$, and when $p_H$ falls
below roughly 0.7 the \hate{} channel no longer wins; two miss by exactly
0.001.

\subsection{The Corrections Fusion Does Not Make}
\label{sec:missed_corrections}

Among the 827 AI errors on ViHSD, 148 non-\clean{} comments have the rule
firing in the correct direction and 151 \clean{} comments have the rule
correctly silent, so the rule agrees with the gold label on 299 errors.  Fusion
acts on 23 of them (7.7\%).  The remaining \textbf{132 missed corrections} are
non-\clean{} cases where the rule already points the right way and the hybrid
is still wrong.  Their median rule \hate{} score is 0.7609 (quartiles 0.67--0.85),
indistinguishable from the 0.7652 median of the corrected cases.  They are
missed because the \clean{} channel keeps the full AI weight ($\alpha_C=1.0$),
so a confidently wrong \clean{} cannot be outvoted: the median fused \clean{}
score is 0.97 for missed cases versus 0.46 for corrected ones.

\subsection{What the Rule Channel Responds To (Cue Taxonomy)}
\label{sec:cue_taxonomy}

We annotated the 16 \hate{} recoveries with a multi-label cue taxonomy: ten
contain teencode or abbreviated profanity (\texttt{đ*o}, \texttt{dm},
\texttt{cmnr}), five standard-spelling profanity or an animal slur, and eight
derogatory wording without profanity (\emph{ngu}, \emph{thua}, \emph{nhục}).
The categories overlap.  The channel is therefore not merely a profanity
detector: much of its value comes from derogatory wording a slur list would
miss, and teencode is its largest single category.

\subsection{Component Ablation (Rule Engine Switches)}
\label{sec:ablation}

Re-tuning and re-freezing the fusion once per switch, only one switch matters:
removing the per-entry confidence costs 0.034 standalone macro-F1 (0.5258 to
0.4914), almost all of it on \off{} (F1 0.2389 to 0.0600), and no other
switch moves the standalone engine by more than 0.008.  Rescaling severity
rewrites 1{,}046 of 6{,}680 rule scores yet changes no hybrid decision
($p=1.0$): $\alpha$ and $\tau$ re-tune.  Dropping the obfuscation views trades
recall for precision (firings 1{,}052 $\to$ 779, precision on fired comments
58.75\% $\to$ 69.58\%); fuzzy matching never fires (0 matches), and case-folding
cannot be ablated because all 6{,}680 test comments are already lower-case.
Inside the hybrid no switch leaves noise ($|\Delta|\le0.0012$): the channel
helps as a channel, not through its weights.

\subsection{Language Detection and Routing Quality}
\label{sec:lang_quality}

Language identification is evaluated across two complementary benchmarks:
(1)~\textbf{Standard sentences} (59{,}868 clean sentences from ViNLI and Leipzig Wortschatz): the detector reaches 99.81\% overall accuracy (99.97\% VI, 99.64\% EN) at 1.67\,ms mean latency (599.5 texts/s; P50 1.63\,ms, P95 1.90\,ms).
(2)~\textbf{Noisy social-media stream} (9{,}158 mixed comments from ViHSD and Davidson): strict routing accuracy is 94.11\% (6{,}658 VI and 2{,}479 EN correctly routed). Short teencode and diacritic drops prompt conservative handling; fallback to \texttt{multilingual} occurs in only 21 cases, preserving robust dual-engine coverage.

\subsection{Posterior Calibration and Blind-Spot Coverage}
\label{sec:calibration_blindspot}

To assess the quality of uncertainty estimates, Table~\ref{tab:hybrid-calibration}
reports calibration diagnostics (Expected Calibration Error, Brier score,
Negative Log-Likelihood, and mean confidence) for AI-only and Hybrid variants on
the ViHSD test split, as well as on the Davidson English test split.
On ViHSD, the Hybrid achieves lower ECE (0.097 vs 0.111, a 12.6\% relative improvement)
and lower Brier score (0.2273 vs 0.2334), indicating that fused probabilities remain well-calibrated
despite the fusion mechanism.  This is important for deployment: a system claiming
high confidence should be trusted, and rule-backed decisions preserve confidence
while uncertain cases benefit from conservative per-channel thresholds.

Among the 504 ViHSD comments missed by all six AI backbones (Table~\ref{tab:blind-spot-504}),
the Rule Engine recovers 124 (24.6\%), split across HATE (39) and CLEAN (85) classes.
The remaining 380 hard cases are errors by both subsystems, indicating fundamental
ambiguity or minority-class sparsity.  The 290 cases where Rule Engine alone succeeds
(and AI fails) confirm that lexical patterns provide orthogonal signal; however, the
low coverage of rule firing (15.75\% of all comments) limits the ceiling on hybrid
improvement.  This trade-off is by design: a conservative system that rarely flags
clean text reduces false positives at the cost of missing some violations.

\begin{table}[t]
\centering
\caption{Calibration diagnostics (ECE, Brier, NLL, mean confidence, Macro-F1) for AI-only vs Hybrid fusion on the bilingual test sets (ViHSD test $n{=}6{,}680$; Davidson EN test $n{=}2{,}478$).}
\label{tab:hybrid-calibration}
\resizebox{\textwidth}{!}{%
\begin{tabular}{llcccccc}
\toprule
Variant & Dataset & $n$ & ECE $\downarrow$ & Brier $\downarrow$ & NLL $\downarrow$ & Mean Conf. $\uparrow$ & Macro-F1 $\uparrow$ \\
\midrule
AI-only & ViHSD test & 6,680 & 0.111 & 0.2334 & 0.8382 & 0.987 & 0.6575 \\
Hybrid & ViHSD test & 6,680 & 0.097 & 0.2273 & 0.7879 & 0.952 & 0.6593 \\
AI-only & EN test (Davidson) & 2,478 & 0.101 & 0.1621 & 0.3212 & 0.808 & 0.7684 \\
Hybrid & EN test (Davidson) & 2,478 & 0.106 & 0.1559 & 0.3285 & 0.804 & 0.7674 \\
\bottomrule
\end{tabular}%
}
\end{table}


% AUTO-GENERATED BY scripts/58_blind_spot_analysis.py - DO NOT EDIT BY HAND
\begin{table}[t]
\centering
\scriptsize
\setlength{\tabcolsep}{4pt}
\caption{Disagreement patterns between AI-only and Rule Engine across the 6,680 ViHSD test samples, highlighting the 504 blind-spot missed by all six AI backbones.}
\label{tab:blind-spot-504}
\begin{tabular}{lcc}
\toprule
Disagreement Pattern & Total Count & Pct. in 504 Blind-Spot \\
\midrule
Both AI and Rule Engine correct & 5,384 & 0.0\% \\
AI correct, Rule Engine wrong   & 469   & 0.0\% \\
AI wrong, Rule Engine correct   & 290   & 24.6\% (124 / 504) \\
Both wrong (hard cases)         & 537   & 75.4\% (380 / 504) \\
\midrule
Total all-miss blind-spot       & 504   & 100.0\% \\
\bottomrule
\end{tabular}
\end{table}


\subsection{Cross-Lingual Domain Robustness (Wortschatz Corpus)}
\label{sec:crosslingual_wortschatz}

To assess domain robustness and avoid false positives on clean out-of-domain text,
we evaluate the English backbone on the Wortschatz Leipzig
2015 English News Corpus, a neutral 30k-sentence benchmark. On 9,879 test sentences,
the classifier predicts 99.89\% CLEAN (7 HATE, 4 OFFENSIVE), confirming the model does
not over-flag neutral text, with throughput sustaining 321.8 texts/sec (3.1 ms per sentence).